% !TEX program = tectonic
\documentclass[11pt,a4paper]{article}
\newcommand{\DocumentKind}{Documentação de negócio}
\newcommand{\DocumentVersion}{1.6}
\usepackage{gemdados}

\begin{document}

\GemCover
  {Documentação de Negócio}
  {Analytics para acompanhamento da aprendizagem}
  {Objetivos, perguntas de negócio, dashboard Looker Studio e limitações de dados.}

\section*{Controle do documento}

\begin{tabularx}{\textwidth}{L{35mm}Y}
\toprule
\textbf{Campo} & \textbf{Valor} \\
\midrule
Versão & 1.6 \\
Data de referência & 24 de setembro de 2026 \\
Público & Coordenação, governança, mentores, analistas e participantes \\
Escopo & Objetivos de negócio, dashboard e entregas validadas do projeto DATACAMP ANALYTICS \\
Fontes de validação & GitHub Projects, repositórios da organização, Looker Studio e nuvem Google \\
\bottomrule
\end{tabularx}

\subsection*{Critério}

O GitHub Project possui 77 cards. Este documento considera como entrega somente
os 55 cards marcados como \textbf{Pronto}. Os demais não são apresentados como
resultado concluído. As entregas e visualizações foram conferidas nos repositórios,
no BigQuery e no Looker Studio em 24 de setembro de 2026.

\clearpage
\tableofcontents
\clearpage

\section{Propósito do projeto}

O projeto transforma relatórios de aprendizagem em informações para
acompanhamento e tomada de decisão. O foco é entender como os membros estudam,
avançam, concluem conteúdos e desenvolvem qualificações, garantindo rastreabilidade
e confiabilidade nos indicadores apresentados.

\section{Termos essenciais}

\begin{longtable}{L{40mm}L{106mm}}
\toprule
\textbf{Termo} & \textbf{Significado} \\
\midrule
\endhead
Analytics & Uso de dados para produzir indicadores, compreender resultados e apoiar decisões. \\
DataCamp & Plataforma de aprendizagem que fornece os relatórios usados pelo projeto. \\
GitHub Project & Quadro de gestão que organiza as entregas em cards e registra sua situação. \\
Governança de dados & Regras e responsabilidades para uso, proteção, qualidade e acesso às informações. \\
Dados pessoais & Informações que identificam uma pessoa, como nome, e-mail ou nome de usuário. \\
Código anônimo & Identificador (\code{user\_id}) que substitui dados pessoais nas análises comuns. \\
Pipeline & Sequência organizada de etapas para receber, proteger, tratar e disponibilizar os dados. \\
BI & Inteligência de negócios: relatórios e painéis para acompanhamento de métricas. \\
Looker Studio & Ferramenta de visualização e construção do dashboard de acompanhamento. \\
Views analíticas & Consultas salvas no BigQuery que centralizam regras de cálculo e métricas. \\
XP & Pontos de experiência obtidos nas atividades de aprendizagem. \\
Homologação & Ambiente usado para validar mudanças antes do uso oficial (\code{gem-dados-lake-stg}). \\
Produção & Ambiente oficial destinado ao uso do projeto (\code{gem-dados-lake-prd}). \\
Squad de Dados & Equipe aprovada responsável pelas atividades de dados do projeto. \\
\bottomrule
\end{longtable}

\section{Squad de Dados}

O projeto organiza a informação sem expor diretamente nome, e-mail ou nome de
usuário nas análises. Dessa forma, a governança pode acompanhar resultados do
grupo e investigar situações autorizadas sem disponibilizar dados pessoais para
o consumo analítico comum. O acesso aos nomes fica restrito à Squad de Dados,
quando necessário e autorizado.

\begin{tabularx}{\textwidth}{L{52mm}Y}
\toprule
\textbf{Função} & \textbf{Integrantes} \\
\midrule
Squad Lead & Kauã Henrique Da Silva \\
Engenharia de Dados - Pipeline & Maria Eduarda Mesquita Sá Barbosa; Isabella Almeida Fernandes \\
Engenharia de Dados - Qualidade & Amanda Böckmann; Juliana Damasceno dos Santos \\
Analytics Engineer & Maria Clara Fagundes da Silva; Larissa Kaynne Botelho Marques \\
Analista de Dados / BI & Jeanne Paula Caputo; Grasiela Costa da Silva \\
\bottomrule
\end{tabularx}

\section{Perguntas de negócio atendidas pelo dashboard}

O dashboard foi construído para responder às perguntas que puderam ser calculadas
de forma consistente a partir das fontes disponibilizadas na camada de staging:

\begin{longtable}{L{52mm}L{94mm}}
\toprule
\textbf{Pergunta de negócio} & \textbf{Decisão e acompanhamento apoiado} \\
\midrule
\endhead
Volume de conclusões & Quantos cursos, projetos, capítulos e assessments foram concluídos pelos mentorados. \\
Taxas de conclusão & Proporção de conclusão das diferentes atividades (cursos, projetos, capítulos e assessments). \\
Distribuição de qualificação & Identificação do nível de proficiência em habilidades e distribuição do status de certificações. \\
Desempenho por trilha & Comparação entre total de cursos e média de cursos concluídos por trilha, e taxa de conclusão por trilha. \\
Variação por segmentos & Comportamento dos indicadores segundo recortes de trilha (\textit{track}), período e nível de habilidade. \\
Entregas por equipe & Acompanhamento consolidado do volume de projetos práticos e certificações concluídas por equipe. \\
\bottomrule
\end{longtable}

\section{Informações que sustentam as análises}

O projeto trabalha com dez conjuntos de relatórios da plataforma de
aprendizagem. Eles foram identificados, classificados e ligados aos temas de
negócio abaixo:

\begin{tabularx}{\textwidth}{L{49mm}Y}
\toprule
\textbf{Assunto} & \textbf{Uso no acompanhamento} \\
\midrule
Membros e licenças & Identificação da população acompanhada e de sua situação de acesso. \\
Resumo geral & Visão consolidada de participação, progresso e atividades concluídas. \\
Certificações & Início, encerramento e resultado das certificações. \\
Catálogo de conteúdo & Relação dos conteúdos e características disponíveis. \\
Atividade em cursos & Início, conclusão, progresso, XP e situação dos cursos. \\
Atividade em programas & Participação e conclusão de programas e trilhas. \\
Atividade em projetos & Participação em projetos práticos e registro de conclusão. \\
Avaliações de habilidades & Tentativas, pontuação e nível de conhecimento reportado. \\
Histórico de equipes & Entrada, saída e movimentação dos membros nas equipes. \\
Tempo de aprendizagem & Tempo dedicado aos diferentes tipos de atividade. \\
\bottomrule
\end{tabularx}

A utilização efetiva no dashboard dependeu da disponibilidade e consistência
dos campos para cada indicador, sendo priorizadas as tabelas de staging
\code{stg\_certifications}, \code{stg\_course\_activity}, \code{stg\_program\_activity},
\code{stg\_skill\_assessment} e \code{stg\_project\_activity}.

\section{Dashboard de Acompanhamento (Looker Studio)}

O dashboard consolida os indicadores em três páginas de visualização,
alimentadas diretamente pela camada staging e views analíticas do BigQuery.

\subsection{Páginas do dashboard}

\begin{longtable}{L{42mm}L{104mm}}
\toprule
\textbf{Página} & \textbf{Indicadores e visualizações apresentados} \\
\midrule
\endhead
Progresso \& Volume & Cards de trilhas concluídas, alunos com trilhas concluídas, cursos concluídos, projetos concluídos e assessments concluídos. Gráfico de atividades concluídas por tipo (exercícios, práticas e capítulos). Taxas de conclusão de cursos, capítulos, projetos e assessments calculadas na view \code{vw\_program\_activity}. \\
\addlinespace
Qualificação & Total de certificações obtidas e assessments realizados. Gráficos de distribuição de desempenho por habilidade e status das certificações. Acompanhamento de projetos concluídos por equipe e certificações concluídas por equipe. \\
\addlinespace
Desempenho & Comparação entre quantidade total de cursos e média de cursos concluídos por trilha. Percentual de taxa de conclusão de cursos por trilha. \\
\bottomrule
\end{longtable}

\subsection{Filtros dinâmicos implementados}

\begin{itemize}
  \item \textbf{Trilha (\textit{Track})}: permite isolar o progresso em trilhas de formação específicas;
  \item \textbf{Período (\textit{Period})}: possibilita a análise temporal delimitada das conclusões e atividades;
  \item \textbf{Nível de habilidade (\textit{Skill Level})}: filtra métricas conforme o nível avaliado nos assessments.
\end{itemize}

\section{Limitações e indicadores não implementados}

Durante a exploração e validação das fontes de dados, algumas métricas inicialmente
previstas em wireframe não puderam ser integradas ao dashboard final. Como princípio
de governança, uma métrica somente foi implementada quando os dados permitiam uma regra
de cálculo consistente, rastreável e auditável, sem recurso a aproximações artificiais.

\subsection{Indicadores previstos no wireframe e descartados}

\begin{longtable}{L{56mm}L{90mm}}
\toprule
\textbf{Indicador previsto} & \textbf{Limitação técnica e de negócio identificada} \\
\midrule
\endhead
Tempo de aprendizado semanal por membro & Dados temporais sem granularidade suficiente para determinar o tempo efetivamente dedicado em cada semana individual. \\
\addlinespace
Volume de conclusão por curso & Dados necessários para relacionar volume de conclusão e demais atributos do curso não estavam disponíveis na estrutura necessária. \\
\addlinespace
Média de projetos práticos concluídos por membro & Não foi possível estabelecer de forma consistente a chave relacional necessária entre os dados de projetos e membros. \\
\addlinespace
Programas ou certificações com maior engajamento & Ausência de dados suficientes para relacionar de forma confiável engajamento, programas e certificações. \\
\addlinespace
Percentual de membros que não iniciaram atividades & Impossibilidade de correlacionar com segurança a totalidade da base de membros com os registros de atividade. \\
\addlinespace
Cursos com alto tempo e baixa conclusão & A granularidade temporal disponível não permitia relacionar adequadamente tempo de aprendizado e conclusão por curso. \\
\addlinespace
Relação entre XP obtido e disponível & A proporção entre XP obtido e XP disponível não foi incorporada à versão final do dashboard. \\
\bottomrule
\end{longtable}

\section{Entregas concluídas no GitHub Project}

\subsection{Entendimento dos dados e das regras}

\begin{itemize}
  \item análise dos dez relatórios e de suas colunas;
  \item identificação de dados pessoais e mapeamento de chaves;
  \item levantamento das relações entre membros, cursos, programas, projetos e avaliações;
  \item criação do dicionário inicial das informações;
  \item definição do escopo do painel a partir dos dados viáveis;
  \item documentação das regras de programas, cursos, certificações, avaliações e XP.
\end{itemize}

\subsection{Recebimento e organização das informações}

\begin{itemize}
  \item reconhecimento automático do tipo de cada relatório recebido;
  \item conversão dos arquivos para formato Parquet;
  \item envio e carga dos dados na camada raw do BigQuery;
  \item validação da quantidade de registros entre a origem e o arquivo processado;
  \item correção dos tipos de datas, números e textos;
  \item atualização sem duplicar registros já existentes.
\end{itemize}

\subsection{Proteção de dados pessoais}

\begin{itemize}
  \item substituição dos identificadores pessoais por código anônimo (\code{user\_id});
  \item separação entre informações analíticas e mapa restrito de identidade;
  \item remoção de nome, e-mail e username das tabelas analíticas e do dashboard;
  \item uso do mesmo código anônimo em todas as tabelas relacionadas.
\end{itemize}

\subsection{Camada analítica e visualização}

\begin{itemize}
  \item limpeza e padronização dos dados na camada staging;
  \item view analítica \code{vw\_program\_activity} no BigQuery;
  \item dashboard Looker Studio com 3 páginas operacionais e filtros dinâmicos;
  \item validação de consistência dos indicadores de volume e taxas de conclusão.
\end{itemize}

\section{Governança e privacidade}

O desenho concluído separa três finalidades:

\begin{enumerate}
  \item \textbf{recebimento}: preserva a informação conforme recebida da origem;
  \item \textbf{tratamento}: corrige tipos, padroniza valores e aplica as regras de proteção;
  \item \textbf{análise e visualização}: disponibiliza dados anonimizados para views e painéis.
\end{enumerate}

O vínculo entre o código anônimo e a identidade do membro fica em área
restrita. O dashboard consome exclusivamente \code{user\_id}, preservando a privacidade
dos participantes.

\section{Validação nos ambientes}

\begin{tabularx}{\textwidth}{L{38mm}Y}
\toprule
\textbf{Ambiente} & \textbf{Situação verificada} \\
\midrule
Homologação & Possui os dez relatórios recebidos, dados tratados na camada staging, view analítica (\code{vw\_program\_activity}) ativa e o dashboard Looker Studio conectado e operacional. \\
Produção & Possui a infraestrutura base, esteiras de automação e serviço de recebimento provisionados via Terraform. Modelos de visualização analítica serão promovidos após ciclo de homologação. \\
\bottomrule
\end{tabularx}

\section{Resumo executivo}

O projeto GEM Dados entregou uma solução integrada de acompanhamento de aprendizagem,
com pipeline automatizado na nuvem Google, proteção total a dados pessoais por
anonimização e dashboard analítico no Looker Studio. O painel viabiliza o acompanhamento
de volumes de conclusão, taxas de avanço e distribuição de qualificações
e certificações por equipe. As limitações inerentes à granularidade das fontes foram
mapeadas com transparência, garantindo indicadores rastreáveis, íntegros e auditáveis.

\end{document}
